\documentclass[11pt,a4paper]{appletmanual}
\usepackage[english]{babel}
\manuallang{en}

\appletname{Non-Linear Budget Constraints}
\appletsub{Seven subsidy schemes bend the budget line in different places. The
question is the same in every one: where does the best affordable bundle land,
and does it land on the fold?}
\appletdir{nonlinear-budget}

\begin{document}
\manualtitlepage{Applet manual · Introductory Microeconomics · Universidad de los Andes\\
Translated from the GeoGebra applet “Max. Utilidad — Rest. Presupuestaria No Lineal”.\\
MIT licensed. It opens in a browser; there is nothing to install.}

\manualtoc
\thispagestyle{empty}
\clearpage
\setcounter{page}{1}

% ==========================================================================
\section{What this is}

A textbook budget constraint is a straight line. Almost no public programme
leaves one straight. A quota of free milk, a discount on the first few units, a
place at a state school that is forfeited the moment you buy a private one:
each bends the frontier somewhere different, and the optimum can land exactly
on the fold.

This applet carries seven schemes, the ability to compare two at once, and an
optimum computed in a way that actually finds folds.

\figher[1.0]{kink}{A non-exclusive in-kind subsidy. The first $x_s$ units are
free, and beyond that the market price applies: the frontier has a flat stretch
and then a kink.}

\begin{amkey}
The comparison that matters is between an in-kind subsidy and the cash transfer
that costs the funder exactly the same. If the quota does not bind, the two
tie. If it binds, cash wins, and the applet measures by how much.
\end{amkey}

\subsection{Who it is for}

The budget-constraint topic, and any public-economics class discussing in-kind
transfers against cash. It assumes indifference curves and the idea of an
optimum; it assumes no calculus.

%% Sections shared by every manual, English edition.
%% Pulled in with \input{common-en}. The only thing that differs between
%% applets is \appletfolder, which is also the folder name and the path on the site.

\section{Opening it}

There is nothing to install. The whole application is one \code{index.html}
file: it downloads no libraries, calls no server, and stores nothing outside
your browser.

\subsection{The three ways}

\begin{description}
\item[Double-click the file.] The quickest route. Find
  \code{\appletfolder/index.html} and open it. The browser shows it straight off the
  disk, with an address beginning \code{file://}. It works offline from the
  first moment.
\item[From the web.] If someone has published it, the address ends in
  \code{/\appletfolder/}. Once the page has loaded you can disconnect: it needs the
  network for nothing else.
\item[Served from a folder.] For a class, putting the folder on any static
  server is enough. With Python installed:
  \begin{quote}\code{python3 -m http.server 8000}\end{quote}
  then \code{http://localhost:8000/\appletfolder/} in the browser.
\end{description}

\begin{amnote}
Saving the page with \key{Ctrl}\,+\,\key{S} (or \key{Cmd}\,+\,\key{S} on a
Mac) leaves a copy of your own that still works offline and that can travel on a
memory stick or go out by email.
\end{amnote}

\subsection{Language and theme}

Two pairs of buttons sit at the top right.

\begin{ctrltable}{Control & What it does}
ES / EN & Switches the whole interface between Spanish and English on the spot.
          Nothing is recomputed, so you can switch mid-explanation without
          losing the state you had. \\
Dark / Light & Switches the theme. Dark is the original design; light exists
          for projectors and for printing. \\
\end{ctrltable}

Both choices are remembered in the browser, so next time it opens the way you
left it.

\subsection{Which browser}

Any reasonably recent one: Chrome, Edge, Firefox or Safari, on a computer,
tablet or phone. On a narrow screen the panels stack one above the other
instead of sitting side by side. No account, no permissions, no extensions.


% ==========================================================================
\section{The seven schemes}

With price $p_x$, income $m$, quota $x_s$ and subsidy rate $s$:

\begin{ctrltable}{No. and name & Frontier}
0 · Simple constraint & The usual line. \\
1 · In-kind subsidy, non-exclusive & Flat to $x_s$, then the market rate: a
  kink. \\
2 · Partial in-kind, non-exclusive & Price $(1-s)p_x$ to $x_s$, then $p_x$: a
  shallower kink. \\
3 · Exclusive subsidy (state schools) & Take $x_s$ free \emph{or} buy at $p_x$:
  a vertical drop. \\
4 · Partial exclusive & A discount on the first $x_s$ units, forfeited entirely
  if you exceed the quota. \\
5 · Cash equivalent (full) & Income $m + p_x x_s$, straight. \\
6 · Cash equivalent (partial) & Income $m + s\,p_x x_s$, straight. \\
\end{ctrltable}

Scheme 1 pairs with 5, and 2 with 6: the same cost to whoever funds the
programme, with no strings attached. Putting one against the other is what the
\ui{Compare with} control is for, and it is the point of the applet.

\begin{amnote}
With $s < 0$ the subsidy turns into a tax of rate $\tau = -s$ on the first
$x_s$ units. The slider reaches $-1$, so the same seven schemes also serve for
talking about selective taxes.
\end{amnote}

\subsection{Kink versus notch}

Schemes 1 and 2 produce a \textbf{kink}: the frontier changes slope but does
not break. Schemes 3 and 4 produce a \textbf{notch}: exceeding the quota
forfeits the whole subsidy, so the frontier falls vertically and there are
bundles no amount of spending reaches.

\fig[1.0]{drop}{An exclusive subsidy. The frontier drops vertically at $x_s$:
one more unit of $x$ costs the whole quota.}

% ==========================================================================
\section{The screen}

\fig[1.0]{interface}{The whole screen.}

\begin{description}
\item[Scheme.] All seven, numbered 0 to 6.
\item[Compare with.] Draws a second scheme behind the first, and switches on
  the \ui{Cash advantage} readout.
\item[Preferences.] Seven families, the last one writable.
\item[Prices and income.] $p_x$, $p_y$, $m$.
\item[The programme.] The quota $x_s$ and the rate $s$.
\item[View.] Heat gradient, substitution panel, axes.
\item[Layers.] What gets drawn.
\end{description}

\subsection{The two panels}

\begin{description}
\item[Plane $(x, y)$.] The utility map under a topographic ramp, everything out
  of reach veiled over, the chosen scheme's frontier, the compared one behind
  it, the indifference curve through the optimum, and the vertices.
\item[Utility along the constraint.] $u(x, f(x))$ as $x$ runs along the
  frontier. Its peak \emph{is} the optimum, so ``find the best affordable
  bundle'' becomes ``find the top of this curve''. On schemes 3 and 4 the curve
  visibly breaks in two: that gap is the forfeited subsidy.
\end{description}

\fig[0.62]{drop-path}{Utility along the constraint under an exclusive scheme.
The gap between the two arcs is exactly the forfeited subsidy.}

\subsection{The dark theme}

\fig[0.95]{dark}{The plane and the substitution panel in the dark theme.}

% ==========================================================================
\section{How the optimum is found}

A first-order condition is the wrong instrument here. On these frontiers the
best bundle is very often at a vertex, where the slope is undefined and no
tangency holds; and on the exclusive schemes it sits at the right-hand end of
the first piece with a strictly worse point immediately to its right.

\begin{enumerate}
\item Every piece of the frontier is scanned densely.
\item The best interior sample on each piece is refined by golden section.
\item \emph{Every vertex is tested explicitly}, from both sides of every join.
\item Whichever wins is the optimum. Where it came from is what classifies it
  as a tangency, a kink, a corner, or the edge of a drop.
\end{enumerate}

\begin{ctrltable}{Kind of optimum & What it means}
Tangency & Interior to a piece: the MRS equals the relative price prevailing
  there. The fold is not biting. \\
Kink & Exactly at the elbow. To the left the constraint is too flat and to the
  right too steep: best to stay at the fold. The consumer takes the quota
  exactly. \\
Corner & On an axis: it pays to spend everything on one good. \\
Edge of the drop & On an exclusive scheme, pinned at the edge. One more unit
  would leave the consumer strictly worse off. \\
\end{ctrltable}

% ==========================================================================
\section{The preferences}

\begin{ctrltable}{Family & Form}
Cobb--Douglas & $x^{a} y^{1-a}$ \\
Perfect substitutes & $a x + (1-a) y$ \\
Perfect complements & $\min(a x, (1-a) y)$ \\
Quasilinear & $8a\ln x + y$ \\
Generalised CES & $\bigl[\tfrac{1}{2a} + \tfrac{1}{2(1-a)}\bigr]^{-1}
   (a x^{\rho} + (1-a) y^{\rho})^{1/\rho}$ \\
X inferior & $\ln(x - \underline{x}) - \tfrac{1}{1-\beta}\ln(\bar y - y)$ \\
Custom & Whatever you type \\
\end{ctrltable}

One weight $a$ runs through all of them: the exponent in Cobb--Douglas, the
relative value in the linear and Leontief cases, the share in the CES.

\begin{amnote}
Two notes on the generalised CES. Its leading factor is written internally as
$2a(1-a)$, which it equals identically: the literal $1/(2a) + 1/(2(1-a))$
divides by zero at both ends of the slider, while $2a(1-a)$ simply goes to zero
there. Being a positive constant it rescales utility without touching
preferences or the optimum. And at $\rho = 0$ the CES is undefined while its
limit is Cobb--Douglas, so a slider resting on zero emits that limit rather
than dividing by it.
\end{amnote}

\fig[1.0]{complements}{Perfect complements over an in-kind subsidy. With
Leontief preferences the optimum is always at the elbow of the indifference
curve, and here it competes with the elbow of the frontier as well.}

% ==========================================================================
\section{The controls}

\begin{ctrltable}{Slider & What it is}
$p_x$, $p_y$ & Prices, bounded below at $0.2$. \\
$m$ & Income. \\
$x_s$ & The programme's quota. It is clamped to what the income can actually
  buy, so the allowance never exceeds the budget. \\
$s$ & The subsidy rate, $-1$ to $1$. Negative is a tax. \\
$a$, $\rho$, $\beta$, $\underline{x}$, $\bar y$ & Parameters of the chosen
  preference family; only the ones that family uses appear. \\
\end{ctrltable}

\subsection{View}

\begin{description}
\item[Heat gradient.] Turn it off and the background goes plain. The
  indifference curves keep their height colour at all times: with the map
  showing it agrees with the ground underneath, and with the map off it is the
  only reading of height left.
\item[Substitution panel.] The utility-along-the-constraint reading is a
  supporting view, so it takes the smaller share of the width and can be
  dismissed entirely, giving the plane the full frame.
\item[Axes.] Fixed by default, set by the $x \le$ and $y \le$ boxes. An axis
  that refits on every slider move renumbers as you drag, and then two settings
  cannot be compared by eye --- which is the whole point of putting a scheme
  and its cash equivalent on one picture. Auto-fit stays available, and a
  warning appears if a frontier runs outside a fixed frame rather than being
  silently cropped.
\end{description}

\subsection{Layers}

\begin{ctrltable}{Layer & What it draws}
Background contours & Level curves of utility. \\
Feasible set & Veils what cannot be afforded. \\
Curve through the optimum & The indifference curve passing through it. \\
Optimum & The chosen point. \\
Vertices & The elbows and edges of the frontier. \\
Compared set & The frontier of the scheme being compared against. \\
Grid & The background rule. \\
\end{ctrltable}

% ==========================================================================
\section{Classroom activities}

\begin{activity}{The usual line}
\amset Scheme \ui{0 · Simple constraint}, opening values.
\amexp $x^\ast = 0.69$, $y^\ast = 2.18$, utility $1.539$, cost $0.00$. The kind
of optimum is \ui{Tangency}: the textbook case.
\par\smallskip
It is the baseline the next six schemes are read against. Write it on the board
before going on.
\end{activity}

\begin{activity}{When the quota does not bite}
\amset Scheme \ui{1 · In-kind subsidy, non-exclusive}, opening values. Set
\ui{Compare with} to \ui{5 · Cash equivalent (full)}.
\amexp $x^\ast = 0.96$, $y^\ast = 3.05$, utility $2.158$, cost $3.50$. The
\ui{Cash advantage} reads $0.000$ and the verdict says ``Cash and in-kind
tie''.
\par\smallskip
Look at why: the optimum lands on the stretch where both frontiers coincide.
The quota $x_s = 0.92$ is smaller than what this consumer was going to buy
anyway, so the programme hands over money and nothing else. An in-kind subsidy
that does not constrain \emph{is} a cash transfer.
\end{activity}

\begin{activity}{When the quota does bite}
\amset As above, but raise $x_s$ to $2.20$.
\amexp $x^\ast = 2.20$ exactly --- the whole quota, not one unit more ---
$y^\ast = 3.11$, utility $2.801$, cost $8.36$. The kind of optimum changes to
\ui{Kink} and the \ui{Cash advantage} rises to $+0.217$.
\par\smallskip
At the same cost of $8.36$, cash reaches a higher indifference curve. Ask why a
government would then pick in-kind. The answers --- paternalism, targeting,
externalities from the subsidised good --- are precisely the ones this model
does not contain, which is why it is worth raising them here.
\end{activity}

\begin{activity}{A notch instead of a kink}
\amset Scheme \ui{3 · Exclusive subsidy (state schools)}, opening values.
\amexp The frontier falls vertically at $x_s$: the \ui{Drop at $x_s$} readout
gives $-1.25$. The kind of optimum is \ui{Edge of the drop}.
\par\smallskip
Switch on the \ui{Substitution panel}. The utility-along-the-constraint curve
breaks into two arcs that do not touch, and the gap between them is exactly the
subsidy forfeited by crossing. Ask what happens to a family that needs barely
one unit more than the quota.
\end{activity}

\begin{activity}{Tax instead of subsidy}
\amset Scheme \ui{2 · Partial in-kind, non-exclusive}. Take $s$ down to $-0.6$.
\amexp The fold reverses: the first $x_s$ units now cost more, and the frontier
is \emph{concave} instead of convex. The optimum moves away from $x$.
\par\smallskip
It is a selective tax on the first few units, and the same picture explains it.
Ask which real tax schemes look like this.
\end{activity}

\begin{activity}{Preferences that always end at the elbow}
\amset \ui{Perfect complements} preferences, scheme \ui{1 · In-kind subsidy}.
\amexp With Leontief the optimum is always at the elbow of the indifference
curve. Move $x_s$: as long as the consumer's elbow is to the right of the
quota, the whole subsidy is taken up and then buying continues.
\par\smallskip
Compare with \ui{Perfect substitutes}, where the optimum jumps from one axis to
the other the moment the relative price crosses the marginal rate of
substitution. The two extreme cases bracket the intuition.
\end{activity}

\begin{activity}{Fixed axes are not a whim}
\amset Any scheme with \ui{Compare with} active. Tick and untick \ui{Fit the
axes automatically} while moving $m$.
\amexp With automatic axes the numbers change at every step and the two
frontiers appear to sit still. With fixed axes you can see which one moves and
by how much.
\par\smallskip
That is why the axes are fixed by default: comparison is the object of the
applet, and a comparison whose scale moves on its own is not one.
\end{activity}

\begin{activity}{Build the case where in-kind wins}
\amset Free rein.
\amexp You cannot, and that is the exercise. Whatever the preferences, cash of
the same cost always reaches an indifference curve at least as high, because
the feasible set under cash contains the one under in-kind.
\par\smallskip
The most that can be achieved is a tie, and the tie happens exactly when the
quota does not bind. Letting the class try for five minutes is worth more than
stating it.
\end{activity}

% ==========================================================================
\section{Expression syntax}

\begin{ctrltable}{Element & What is allowed}
Operators & \code{+ - * / \^{} ( )} \\
Functions & \code{sqrt ln log log10 exp abs sin cos tan min max} \\
Constants & \code{e}, \code{pi} \\
Implicit multiplication & \code{2xy} is \code{2*x*y} \\
Associativity & \code{\^{}} binds to the right and tighter than unary minus:
  \code{-x\^{}2} is $-(x^2)$ \\
\end{ctrltable}

% ==========================================================================
\section{Changes from the GeoGebra original}

\begin{description}
\item[Discontinuities are modelled, not approximated.] The original wrote
  schemes 3 and 4 as single \code{If[...]} functions, which GeoGebra draws with
  a connecting segment across the jump. Each scheme here is a list of separate
  pieces, so the drop stays a drop and the utility profile breaks where it
  should.
\item[Colour range.] \code{uMIN := U(0.001, 0.001)} is undefined for the
  Stone-Geary utility, whose domain needs $x > \underline{x}$; the ramp
  collapsed to a single colour for the whole ``x inferior'' example. The
  rewrite takes the 2nd percentile of the sampled field.
\item[\code{uMAX}] used the same x-derived value for both coordinates, and
  exceeded $\bar y$ often enough to return \code{NaN}. Same fix.
\item[Prices.] $p_X$ and $p_Y$ sliders started at 0, making $I/p_Y$ and the
  whole construction divide by zero. Both are bounded below by $0.2$.
\item[$x_s$ and $\underline{x}$] are clamped to what the income can actually
  buy, so the allowance can never exceed the budget and the Stone-Geary utility
  stays defined.
\item[Dead objects.] \code{pXlow} and \code{pXhigh} were computed from a
  formula that returns large negative numbers for every setting the sliders
  allow, and were displayed nowhere. Not carried over.
\item[The cash comparison is computed], not left for the reader to eyeball: the
  cost of each scheme and the utility difference against the equal-cost cash
  transfer are readouts.
\end{description}

% ==========================================================================
\section{Troubleshooting}

\begin{ctrltable}{Symptom & What is happening}
A warning that the frontier runs off & The axes are fixed and the drawing does
  not fit. Raise $x \le$ or $y \le$, or tick auto-fit. \\
The plane is plain & The \ui{Heat gradient} is off. The curves keep their
  height colour regardless. \\
Cash advantage reads ``---'' & There is no scheme in \ui{Compare with}. \\
The optimum ignores $s$ & The quota is not binding: the optimum falls outside
  the subsidised stretch. Raise $x_s$. \\
Utility comes out \code{NaN} & In the ``X inferior'' family the domain needs
  $x > \underline{x}$ and $y < \bar y$. Adjust those two sliders. \\
\end{ctrltable}

%% Sections shared by every manual, English edition.
%% Pulled in with \input{common-en}. The only thing that differs between
%% applets is \appletfolder, which is also the folder name and the path on the site.

\section{Publishing it for a class}

Any static host will do: GitHub Pages, Netlify, a university web directory,
even a shared folder served over HTTP. The only thing to upload is the
\code{index.html} file, inside a folder named however you want the address to
read.

The repository ships a GitHub Actions workflow
(\code{.github/workflows/pages.yml}) that publishes every applet at once on each
push to the default branch. To switch it on, in the repository:
\emph{Settings\,$\rightarrow$\,Pages\,$\rightarrow$\,Source: GitHub Actions}.
That is the one step the workflow cannot do for itself, because creating a Pages
site needs administration rights that a \code{GITHUB\_TOKEN} is never granted.

\begin{amnote}
To hand the applet to students without reliable internet, send them the
\code{index.html} file directly. It weighs less than a photograph and opens with
a double-click.
\end{amnote}

\section{Licence and provenance}

The application code is MIT. Use it, change it and pass it on, in class or out
of it, with attribution.

This applet is a standalone rewrite of an original GeoGebra applet. The rewrite
is not a literal translation: where the original had a construction error, a
division by zero reachable with its own sliders, or a dead object inherited from
whatever applet it was copied from, this version fixes it and says so. The
section \emph{Changes from the GeoGebra original} lists those changes one by
one.

Everything is hand-written and dependency-free: the expression parser, the
symbolic differentiation, the level-curve tracing and the drawing. In
particular there is no \code{eval}, so what a student types into a text box
cannot become code, and a strict content security policy does not break the
page.


\end{document}
